本研究は, 状態空間モデルを用いて日本の宅配便需要に与えた外部イベント（運賃改定, コロナ禍）の影響を分離・定量化することを目的とする. 
モデル比較の結果, コロナ禍の影響を複数フェーズで捉えたモデルが最も優れており, 
そのインパクトが持続的な構造変化と短期的な需要急増の複合現象であることを明らかにした. 
本稿は, 状態空間モデルが複雑なイベント効果の分離に有効であることを実証し, 物流の精緻な把握に資するものである. 